\begin{textsample}{NEG Dim 2 – 2011\_c5 – Score -2.00 – 2011\_c5/enade\_2011\_ef\_5f734ad795.txt}  \label{ex:f2_neg_090}
CAPÍTULO 5 DISTRIBUIÇÃO DOS CONCEITOS 5.1 PANORAMA NACIONAL DA DISTRIBUIÇÃO DOS CONCEITOS A Tabela 5.1 apresenta a quantidade e distribuição de cursos de Educação Física participantes do ENADE/2011, por faixa de conceito e Grande Região.
A diferença entre os cursos tabulados neste capítulo e no capítulo 2 são os cursos sem conceito, em princípio, aqueles sem alunos concluintes que participassem da prova.
Observando -se os dados da Tabela 5.1, nota -se que, dos 466 cursos participantes, 225 ( 48,3 \% ) classificaram -se com conceito 3, o valor modal.
Este foi também o conceito modal em todas as regiões, com proporções variando de 42 \% a 54 \%.
O conceito 2 foi o segundo mais frequente em nível nacional ( 21,0 \%, correspondendo a 98 cursos ) e o conceito 4, o terceiro ( 20,4 \%, correspondendo a 95 cursos ).
Houve, ainda, 17 cursos ( 3,6 \% ) que receberam conceito 5 e outros 25 cursos que receberam conceito 1.
Dentre os cursos de Educação Física seis cursos ficaram sem conceito ( SC ).
A região Norte participou com 26 cursos ou 5,6 \% do total nacional.
Nesta região, o conceito 3 foi atribuído a 14 cursos dos 26 participantes, o que equivale a 53,9 \% do total regional.
Nenhum curso ficou sem conceito ou recebeu o conceito 1.
Os demais foram avaliados com conceito 2 ( cinco cursos correspondendo a 19,2 \% ), conceito 4 ( seis cursos correspondendo a 23,1 \% ) e conceito 5 ( um curso correspondendo a 3,8 \% ).
A região Nordeste participou com 61 cursos ou 13,1 \% do total nacional.
Destes, 29 cursos, 47,6 \% em termos regionais, obtiveram conceito 3, o conceito modal para a região.
Os conceitos 2 e 4 foram atribuídos a, respectivamente, dez e 13 cursos ( 16,4 \% e 21,3 \% ).
O conceito 5 foi atribuído a um curso ( 1,6 \% ) e o conceito 1 a sete cursos ( 11,5 \% ).
Nessa região um dos cursos ficou sem conceito.
Dos 229 cursos participantes da região Sudeste ( 49,1 \% do total nacional ), o conceito modal foi 3, atribuído a 101 ( 44,1 \% ) dos cursos da região.
O conceito 1 foi atribuído a 11 cursos ( 4,8 \% ) e o conceito 2 a 57 ( 24,9 \% ).
O conceito 4 foi alocado a 49 cursos ( 21,4 \% ) e oito outros ( 3,5 \% ) receberam o conceito 5.
Três dos cursos ficaram sem conceito.
A região Sul também contou com cursos \textbf{distribuídos} em todas as faixas de conceitos.
A predominância do conceito 3 foi de 58,3 \%, correspondentes a 63 dos 108 cursos participantes na região Sul.
O conceito 2 foi atribuído a 12 cursos ( 11,1 \% ) e o conceito 4 a 23 cursos ( 21,3 \% ).
O conceito 1 foi recebido por cinco cursos ( 4,6 \% ) e o conceito 5, por três deles.
Dois dos cursos da região Sul ficaram sem conceito.
Quase metade ( 42,9 \%, correspondendo a 18 cursos ) dos 42 cursos participantes na região Centro-Oeste recebeu conceito 3, o conceito modal.
Os demais cursos foram avaliados com : conceito 2 ( 14 cursos, 33,3 \% ), conceito 4 ( quatro cursos, 9,5 \% ) e conceito 5 ( outros quatro cursos, 9,5 \% ).
Nesta região nenhum curso ficou sem conceito. 5.2 CONCEITOS POR CATEGORIA ADMINISTRATIVA E POR GRANDE REGIÃO A Tabela 5.2 apresenta a distribuição dos cursos participantes do ENADE/2011 de Educação Física, por Categoria Administrativa, de acordo com os conceitos por eles alcançados, segundo as Grandes Regiões brasileiras.
Dos 466 cursos participantes, 106 ( 22,7 \% ) eram ministrados em instituições públicas e 360 ( 77,3 \% ), em privadas.
De acordo com as informações da Tabela 5.2, em termos proporcionais houve equilíbrio na obtenção do conceito 5 entre instituições públicas e privadas : das 17 IES que receberam este conceito, 4 eram públicas e 13 privadas, correspondendo a 3,8 \% e 3,6 \% das respectivas categorias.
Dos 106 cursos participantes de IES públicas, o conceito 3 foi o valor modal, atribuído a 35 cursos ( 33,0 \% ).
Entre os demais cursos participantes desta categoria administrativa, 18 obtiveram conceito 1 ( 17,0 \% da categoria ), 17 receberam conceito 2 ( 16,0 \% ) e 30 foram avaliados com conceito 4 ( 28,3 \% ).
Dentre os demais, quatro receberam o conceito 5 e dois dos cursos ficaram sem conceito.
Na rede privada, o conceito modal também foi 3, com 190 cursos dos 360 da categoria ( 52,8 \% ).
Entre os demais cursos privados participantes, sete receberam conceito 1 e 81 o conceito 2.
O conceito 4 foi atribuído a 65 cursos e o conceito 5 a 13 cursos.
Nesta categoria administrativa, quatro dos cursos ficaram sem conceito.
Na análise por região, observa -se que na região Norte, as instituições privadas participaram com 15 cursos ( 57,7 \% do total regional ), dos quais nenhum ficou sem conceito.
O conceito modal para as instituições privadas na região foi o 3, com nove cursos, correspondendo a 60,0 \%.
Quatro cursos receberam o conceito 2 ( 26,7 \% ) e dois outros cursos receberam conceito 4 ( 13,3 \% ).
Nenhum curso foi avaliado com conceitos 1 e 5.
As instituições públicas participaram com 11 cursos ( 42,3 \% do total regional ), dos quais nenhum ficou sem conceito.
O conceito modal para esta Categoria Administrativa, na região Norte, foi o 3, com cinco cursos, correspondendo a 45,5 \%.
Os demais cursos receberam os conceitos : 2 ( um curso ), 4 ( quatro cursos ) e 5 ( um outro curso ).
Nenhum curso obteve conceito 1 nesta categoria na região.
Na região Nordeste, a rede privada concentrou 32 dos 61 cursos participantes, equivalentes a 52,5 \% do total da região, a menor proporção de cursos em rede privada entre todas as regiões.
Nesta Categoria Administrativa na região, \textbf{predominaram} os cursos com conceito 3, 15 cursos correspondendo a 46,9 \%.
Os demais foram avaliados com conceitos : 1 ( três cursos correspondendo a 9,4 \% ), 2 ( dez cursos correspondendo a 31,3 \% ) e 4 ( quatro cursos correspondendo a 12,5 \% ).
Nenhum curso recebeu o conceito 5 ou ficou sem conceito.
As instituições públicas dessa região participaram com 29 cursos ( 47,5 \% dos cursos da região ), dos quais 14 obtiveram conceito 3, o valor modal.
Quatro cursos obtiveram conceito 1, nove cursos obtiveram conceito 4 e um curso obteve o conceito 5.
Um curso ficou sem conceito e nenhum curso recebeu conceito 2.
Na região Sudeste, a proporção de cursos da rede privada, 86,9 \%, a proporção mais elevada dentre as regiões brasileiras, correspondendo a 199 dos 229 cursos participantes.
Nesta categoria, na região Sudeste, o conceito modal foi 3 ( 96 cursos ).
Os demais foram avaliados com conceito 1 ( quatro cursos ), conceito 2 ( 48 cursos ), conceito 4 ( 42 cursos ), e conceito 5 ( sete cursos ).
Nesta combinação de Categoria Administrativa e Grande Região, dois cursos não receberam conceito.
Entre os 30 cursos oferecidos em instituições públicas na região Sudeste, a categoria modal foi a que recebeu conceito 2 ( nove cursos ).
Os demais foram avaliados com os conceitos 1 ( sete cursos ), 3 ( cinco cursos ), 4 ( sete cursos ) e conceito 5 ( um curso ).
Na região Sudeste nesta categoria, um curso ficou sem conceito.
As instituições privadas concentraram 87 dos 108 cursos participantes da região Sul, 80,6 \% do total regional.
Desses, 57 ficaram com conceito 3, o conceito modal.
Os demais foram avaliados com conceito 2 ( dez cursos ), conceito 4 ( 15 cursos ) e conceito 5 ( três cursos ).
Nesta combinação de Categoria Administrativa e Grande Região, dois cursos ficaram sem receber conceito e nenhum curso recebeu conceito 1.
As instituições públicas na região Sul participaram com 21 cursos ( 19,4 \% ), dos quais nenhum ficou sem conceito.
O conceito modal foi 3, com seis cursos.
O conceito 1 foi recebido por cinco cursos, o 2 por dois cursos e o 4 por oito cursos.
Nenhum curso alcançou o conceito 5.
Na região Centro-Oeste, 27 dos 42 cursos participantes eram de instituições privadas ( 64,3 \% em termos regionais ).
Destes, 13 concentraram -se no conceito 3, conceito modal.
Os demais receberam conceito 2 ( nove cursos ), conceito 4 ( dois cursos ) e conceito 5 ( três cursos ).
Nenhum curso ficou sem conceito nesta categoria, nem tampouco foi recebeu conceito 1.
Dos 15 cursos de instituições públicas, cinco obtiveram conceito 2 e também cinco receberam 3, conceitos modais.
Os conceitos 1 e 4 foram alocados a dois cursos cada e o conceito 5 foi alocado a um curso.
Nesta região, nenhum curso de IES pública ficou sem conceito. 5.3 CONCEITOS POR ORGANIZAÇÃO ACADÊMICA E POR GRANDE REGIÃO Na Tabela 5.3 encontra -se a distribuição dos conceitos atribuídos aos cursos participantes do ENADE/2011 na Área de Educação Física, por Organização Acadêmica, segundo as Grandes Regiões brasileiras.
Dos 466 cursos participantes, 220 eram oferecidos em Universidades, 80 em Centros Universitários e os demais 166 em Faculdades.
Esta distribuição corresponde a, respectivamente, 47,2 \%, 17,2 \% e 35,6 \% dos cursos.
De acordo com os dados apresentados, 15 dos 17 cursos avaliados com conceito 5 eram vinculados a Universidades.
Cinco cursos em Universidades ficaram sem conceito e esse tipo de Organização Acadêmica teve o conceito 3 como modal, com 101 cursos.
Os demais cursos avaliados receberam os conceitos : 1 ( 20 cursos ), 2 ( 24 cursos ), 4 ( 55 cursos ) e conceito 5 ( 15 cursos, como já mencionado ).
Entre os cursos em Centros Universitários, o conceito modal, também foi o 3, com 47 cursos.
Neste tipo de organização acadêmica nenhum curso ficou sem conceito.
Os outros cursos neste tipo de Organização Acadêmica receberam os conceitos : 1 ( dois cursos ), 2 ( 15 cursos ), 4 ( outros 15 cursos ) e 5 ( um curso ).
Nas Faculdades, um dos 166 cursos ficou sem conceito e 77 receberam o conceito modal 3.
Dos demais cursos neste tipo de Organização Acadêmica, três receberam conceito 1, 59 receberam o conceito 2, 25 receberam o conceito 4 e um recebeu o conceito 1.
Considerando -se separadamente as regiões brasileiras, verifica -se que, na região Norte nenhum curso ficou sem conceito ou recebeu o conceito 1.
As Universidades concentraram 11 dos 26 cursos participantes.
Dentre estes, seis cursos receberam o conceito 3, conceito modal, quatro receberam o conceito 4 e um no conceito 5.
Os Centros Universitários da região Norte foram representados por sete cursos.
Três receberam o conceito 3 e os quatro restantes dividiram -se igualmente entre os conceitos 2 e 4.
As Faculdades participaram com oito cursos na região Norte, cinco receberam o conceito 3, conceito modal, e os outros três o conceito 2.
Na região Nordeste, as Universidades participaram com 31 dos 61 cursos na Área de Educação Física da região.
Um dos cursos oferecidos em Universidades do Nordeste ficou sem conceito.
O conceito modal foi 3, com 14 cursos.
Os demais receberam conceito 1 ( cinco cursos ), conceito 2 ( um curso ), 4 ( nove cursos ) e conceito 5 ( também um curso ).
Os Centros Universitários contaram com três cursos participantes na região Nordeste, sendo dois com conceito 3 e um com conceito 2.
As Faculdades foram representadas por 27 cursos na região Nordeste e nenhum ficou sem conceito.
O conceito modal foi 3, com 13 cursos.
Dos cursos restantes, dois receberam conceito 1, oito receberam conceito 2 e quatro receberam conceito 4.
Na região Sudeste, as Universidades concentraram 85 dos 229 cursos da região.
Entre os cursos em Universidades, o conceito modal foi 3 com 33 cursos, e dois cursos ficaram sem conceito.
Os demais cursos de Educação Física oferecidos em Universidades no Sudeste receberam os conceitos : 1 ( oito cursos ), 2 ( 14 cursos ), 4 ( 21 cursos ) e 5 ( sete cursos ).
Os Centros Universitários participaram com 53 cursos na região Sudeste, dos quais 34 obtiveram conceito modal, 3, e nenhum ficou sem conceito.
Os demais receberam os conceitos 1 ( dois cursos ), 2 ( seis cursos ) e conceito 4 ( 11 cursos ).
Nesta combinação de organização acadêmica e região nenhum curso alcançou conceito 5.
As Faculdades foram representadas por 91 cursos na região Sudeste, que se \textbf{distribuíram} nos conceitos : 1 ( um curso ), 2 ( 37 cursos, conceito modal ), 3 ( 34 cursos ), 4 ( 17 cursos ) e 5 ( também um curso ).
Um curso ficou sem conceito.
Dos 108 cursos da região Sul, 71 eram de Universidades, para os quais o conceito modal foi 3, com 39 cursos.
Nesse tipo de organização, dois dos cursos ficaram sem conceito e os demais receberam os conceitos : 1 ( cinco cursos ), 2 ( quatro cursos ), 4 ( 18 cursos ) e 5 ( três cursos ).
Os Centros Universitários da região Sul tiveram cinco dos 11 cursos participantes no conceito modal, 3.
Os outros cursos receberam conceitos 2 ( quatro cursos ) e 4 ( dois cursos ).
Foram 26 os cursos vinculados a Faculdades na região Sul e nenhum deles ficou sem conceito.
Os demais receberam conceito 2 ( quatro cursos ), conceito 3 ( 19 cursos, o conceito modal ) e conceito 4 ( três cursos ).
Na região Centro-Oeste, 22 dos 42 cursos eram de Universidades.
Nesse tipo de organização, nenhum curso ficou sem conceito e o conceito modal foi 3, com nove cursos.
Os outros cursos obtiveram os conceitos 1 ( dois cursos ), 2 ( cinco cursos ), 4 ( três cursos ) e 5 ( outros três cursos ).
Os Centros Universitários da região Centro-Oeste contaram com seis cursos, dos quais três receberam o conceito 3 ( conceito modal ), dois receberam conceito 2 e um curso o conceito 5.
Nenhum curso ficou sem conceito.
Dos 14 cursos em Faculdades na região Centro-Oeste, nenhum ficou sem conceito, sete receberam conceito 2, seis o conceito 3 e um o conceito 4.

% matched lemmas: distribuir, predominar
\end{textsample}
